\section{Minimax Lower Bound}

This section gives the observed-law minimax lower bound over the class \(\mathcal P_{\alpha,\gamma}\) in \cref{obj:def:law-class}.  The argument uses a two-point testing construction in the sense of \citet{LeCam1986}: two laws agree away from a small active block, differ only in the sign of a local treatment contrast on that block, and are calibrated so that their product distributions remain statistically close.  The regret separation is then read through the welfare identity in \cref{obj:thm:welfare-identity}.  The result characterizes the observed-law converse exponent for this law class.

\begin{definitionv}{def:minimax-regret}[Minimax regret $M_n(\alpha,\gamma)$]
The minimax regret over the law class \(\mathcal P_{\alpha,\gamma}\) is
\[
M_n(\alpha,\gamma)
=
\inf_{\widehat\pi_n}
\sup_{P\in\mathcal P_{\alpha,\gamma}}
\mathbb E_P\!\left[
R_P(\widehat\pi_n)
\right],
\]
where the infimum ranges over all measurable data-dependent estimators \(\widehat\pi_n\) taking values in \(\Pi\).
\end{definitionv}

The risk in \cref{obj:def:minimax-regret} is evaluated under the observed law \(P\), and the supremum is taken only over laws satisfying the boundedness, positivity, margin, zero-effect, overlap-decay, and endpoint restrictions collected in \cref{obj:def:law-class}.  The estimator is required to output an element of the policy class \(\Pi\), while regret remains measured relative to the oracle rule in \cref{obj:def:welfare-regret}.

The lower bound is driven by the balance between three quantities: the block mass allowed by the margin condition, the local contrast size, and the treatment-arm sampling probability on the weak-overlap side.  The construction below fixes the contrast scale \(h_n\), makes the active block have mass of order \(h_n^\alpha\), and assigns the weak treatment arm probability of order \(h_n^{\beta_{\alpha,\gamma}}\) when overlap decay permits it.

\begin{definitionv}{def:two-point-witness}[Two-point witness laws \(P_{n,\sigma}\)]
Let \(\mathcal X=[0,1]\) and let \(P_X\) be Lebesgue measure restricted to \([0,1]\). Set
\[
h_n=n^{-1/(2+\alpha+\beta_{\alpha,\gamma})},
\qquad
q_n=
\begin{cases}
1/4, & \beta_{\alpha,\gamma}=0,\\
h_n^{\beta_{\alpha,\gamma}}, & \beta_{\alpha,\gamma}>0,
\end{cases}
\]
and
\[
B_n=\{x\in\mathbb R:0\le x\le c_Bh_n^\alpha\}.
\]
Set \(\tau_0=(u_0+2)/2\). For each \(\sigma\in\{1,-1\}\), the law \(P_{n,\sigma}\) has the same covariate law \(P_X\), the same logging propensity
\[
e(x)=
\begin{cases}
q_n, & x\in B_n,\\
1/2, & x\notin B_n,
\end{cases}
\]
and the same off-block outcome law supported on \(\{-1,1\}\), with
\[
E[Y\mid X=x,A=1]=\tau_0/2,
\qquad
E[Y\mid X=x,A=0]=-\tau_0/2,
\qquad x\notin B_n .
\]
Thus the common off-block optimal label is \(1\), and the common off-block contrast is \(\tau_0\). On \(B_n\), set \(Y=0\) almost surely when \(A=0\). On the active treated cell \(\{X\in B_n,A=1\}\), let \(Y\) be supported on \(\{-1,1\}\) with
\[
P(Y=1\mid X\in B_n,A=1)=\frac{1+\sigma h_n}{2},
\qquad
P(Y=-1\mid X\in B_n,A=1)=\frac{1-\sigma h_n}{2}.
\]
Consequently,
\[
E[Y\mid X\in B_n,A=1]=\sigma h_n,
\qquad
\tau_{P_{n,\sigma}}(x)=\sigma h_n
\quad\text{for }x\in B_n .
\]
The two laws agree everywhere except on the treated active cell. The constants satisfy
\[
8c_B<\log 5 .
\]
\end{definitionv}

The construction in \cref{obj:def:two-point-witness} is one-sided in the relevant sense: the two laws differ only in the treated observations on \(B_n\), and the probability of entering that informative cell is \(P_X(B_n)q_n\).  Thus weak overlap reduces information exactly where the sign of the contrast is hardest to learn.

The exponent \(\beta_{\alpha,\gamma}\) in \cref{obj:def:exponents} is the largest weak-arm exponent compatible with the overlap-decay restriction on a block of margin mass \(h^\alpha\).  The next lemma records this calibration.

\begin{lemmav}{prop:overlap-envelope}[Overlap envelope calibration]
For \(\alpha\ge0\), \(\gamma>0\), \(h\in(0,1)\), and \(\beta\ge0\), consider the tight window
\[
v=h^\beta,
\qquad
u=h^{\beta/\gamma}.
\]
Then
\[
u^\alpha v^{1/\gamma}
=
h^{(\alpha+1)\beta/\gamma},
\]
and
\[
u^\alpha v^{1/\gamma}\ge h^\alpha
\quad\Longleftrightarrow\quad
\beta\le \beta_{\alpha,\gamma}=\frac{\alpha\gamma}{\alpha+1},
\]
with equality in the displayed power comparison whenever \(\beta=\beta_{\alpha,\gamma}\). Moreover \(\beta_{\alpha,\gamma}\ge0\), and for every \(\beta'\ge0\), the corresponding tight window \(v=h^{\beta'}\), \(u=h^{\beta'/\gamma}\) satisfies
\[
u^\alpha v^{1/\gamma}\ge h^\alpha
\quad\Longleftrightarrow\quad
\beta'\le\beta_{\alpha,\gamma}.
\]
Thus, within this algebraic tight-window envelope, \(\beta_{\alpha,\gamma}\) is the largest admissible weak-arm exponent.
\end{lemmav}

\cref{obj:prop:overlap-envelope} is the source of the denominator \(2+\alpha+\beta_{\alpha,\gamma}\).  The term \(2\) comes from distinguishing Bernoulli means separated by \(h_n\), the term \(\alpha\) from the block mass allowed by the margin condition, and the term \(\beta_{\alpha,\gamma}\) from the probability of observing the informative treatment arm.

The witness laws must also satisfy the restrictions defining \(\mathcal P_{\alpha,\gamma}\).  The next statement verifies membership and records the two oracle rules induced by the sign change on \(B_n\).

\begin{lemmav}{lem:witness-membership}[Witness-law membership]
Assume \(\alpha\ge0\), \(\gamma\ge0\), the margin-window normalization \(0<u_0<2\) in \cref{obj:ass:margin-window}, \(C_m,C_o,c_o,c_B,\underline p>0\), \(c_B\le C_m\), \(c_B\le C_o\), and \(\underline p\le 1/4\). Also assume the witness constant satisfies the overlap-decay smallness conditions
\[
0<\gamma,\ 0<\alpha
\quad\Longrightarrow\quad
c_B\le C_o c_o^{-\alpha/\gamma},
\qquad
0<\gamma,\ \alpha=0
\quad\Longrightarrow\quad
c_B\le C_o 4^{-1/\gamma}.
\]
Then, for each witness sign \(\sigma\in\{+,-\}\), for all sufficiently large \(n\), the witness law \(P_{n,\sigma}\) of \cref{obj:def:two-point-witness} satisfies the observed-law membership conditions defining the class \(\mathcal P_{\alpha,\gamma}\) in \cref{obj:def:law-class}. Hence, for all sufficiently large \(n\),
\[
P_{n,+},P_{n,-}\in \mathcal P_{\alpha,\gamma}.
\]
Moreover, the two corresponding witness-optimal policies are
\[
x\mapsto 1
\qquad\text{and}\qquad
x\mapsto \mathbf 1\{x\notin B_n\},
\]
where \(B_n=\{x:0\le x\le c_B h_n^\alpha\}\). If these two policies belong to \(\Pi\), then they are available as the two policy actions in the associated two-point lower-bound reduction. The construction uses the off-block contrast \(\tau_0=(u_0+2)/2\), so \(\tau_0\in(u_0,2)\) under the normalization \(0<u_0<2\).
\end{lemmav}

\cref{obj:lem:witness-membership} has two roles.  First, it places both local alternatives inside the same observed-law class used in the minimax risk.  Second, it identifies the two witness-optimal rules.  When those two rules lie in \(\Pi\), they are the two admissible policy actions in the associated two-action testing reduction; the regret lower bound itself is still stated for the minimax risk in \cref{obj:def:minimax-regret}, with regret evaluated relative to the oracle rule.

It remains to check that the two laws are hard to distinguish at sample size \(n\).  Because the laws differ only on a set with covariate mass \(m_n\), treatment probability \(q_n\), and conditional mean separation of order \(h_n\), the one-observation divergence has order \(m_nq_nh_n^2\).

\begin{lemmav}{lem:two-point-divergence}[Two-point divergence]
For the witness laws in \cref{obj:def:two-point-witness}, under \cref{obj:ass:margin-window}, assume in addition that
\[
0<c_B,\qquad \alpha\ge 0,\qquad \gamma\ge 0,
\qquad 8c_B<\log 5.
\]
Then there exists a constant \(C>0\) such that, for all sufficiently large \(n\),
\[
P_{n,+}^{\otimes n}\ll P_{n,-}^{\otimes n},
\qquad
\int \left(\frac{dP_{n,+}^{\otimes n}}{dP_{n,-}^{\otimes n}}-1\right)^2
  dP_{n,-}^{\otimes n}<\infty.
\]
Moreover the one-observation chi-square divergence satisfies
\[
\chi^2(P_{n,+}\,\|\,P_{n,-})
 \le C\,h_n^{2+\alpha+\beta_{\alpha,\gamma}},
\qquad
h_n=n^{-1/(2+\alpha+\beta_{\alpha,\gamma})}.
\]
For the same constant \(C\), the product divergence is bounded:
\[
\chi^2(P_{n,+}^{\otimes n}\,\|\,P_{n,-}^{\otimes n})\le C.
\]
\end{lemmav}

The bandwidth \(h_n\) in \cref{obj:def:two-point-witness} is chosen so that \(n h_n^{2+\alpha+\beta_{\alpha,\gamma}}\) is bounded.  Hence no estimator can reliably determine which sign generated the active block merely from the observed sample.

The same active block creates the welfare separation.  Under one law the oracle treats \(B_n\), and under the other law the oracle does not.  Any single policy must disagree with at least one of these two assignments on a nontrivial part of the block.

\begin{lemmav}{lem:regret-separation}[Regret separation]
For the witness laws in \cref{obj:def:two-point-witness}, suppose \cref{obj:ass:margin-window} holds, \(\alpha,\gamma\ge0\), \(c_B,C_m,C_o,c_o,\underline p>0\), \(c_B\le C_m\), \(c_B\le C_o\), \(\underline p\le 1/4\), every \(\pi\in\Pi\) is measurable, and the block-size constant satisfies
\[
\gamma>0,\ \alpha>0 \implies c_B\le C_o c_o^{-\alpha/\gamma},
\qquad
\gamma>0,\ \alpha=0 \implies c_B\le C_o 4^{-1/\gamma}.
\]
Then there exists a constant \(c>0\) such that, for all sufficiently large \(n\),
\[
\forall\pi\in\Pi,\qquad
\max_{\sigma\in\{+,-\}}
R_{P_{n,\sigma}}(\pi)
\ge
c\,h_n^{1+\alpha}.
\]
For such \(n\), the two optimal policies disagree on the active block. Hence every \(\pi\in\Pi\) misclassifies the active block under at least one of the two laws, and the sum of the two regrets is bounded below by the contrast scale times a covariate subblock of mass \((\min\{c_B,1\})h_n^\alpha\). Passing from the sum to the maximum gives the stated separation of order \(h_n^{1+\alpha}\).
\end{lemmav}

\cref{obj:lem:regret-separation} converts the testing ambiguity into welfare regret.  The separation is \(h_n\) per misclassified point on a set of mass \(h_n^\alpha\), giving \(h_n^{1+\alpha}\).

The testing step uses the standard chi-square form of Le Cam's two-point method \citep{LeCam1986,Tsybakov2009}.

\begin{lemmav}{lem:le-cam-two-point-chisq}[Chi-square testing bound]
If
\[
\chi^2(P_+^{\otimes n}\,\|\,P_-^{\otimes n})\le C_\chi,
\]
then every test \(\psi\) taking values in \(\{+,-\}\) satisfies
\[
P_+^{\otimes n}(\psi\ne +)
+
P_-^{\otimes n}(\psi\ne -)
\ge
c(C_\chi)>0.
\]
In particular, this conclusion holds whenever the one-observation chi-square divergence is bounded by \(c_0/n\), since
\[
1+\chi^2(P_+^{\otimes n}\,\|\,P_-^{\otimes n})
=
\left(1+\chi^2(P_+\,\|\,P_-)\right)^n
\le
\left(1+\frac{c_0}{n}\right)^n
\le e^{c_0}.
\]
\end{lemmav}

Combining \cref{obj:lem:two-point-divergence}, \cref{obj:lem:regret-separation}, and \cref{obj:lem:le-cam-two-point-chisq} yields the minimax lower bound.  Since \(h_n^{1+\alpha}=n^{-(1+\alpha)/(2+\alpha+\beta_{\alpha,\gamma})}\), the lower-bound exponent is \(r_\star(\alpha,\gamma)\).

\begin{theoremv}{thm:minimax-lower}[Minimax lower bound]
Under \cref{obj:ass:margin-window}, suppose that \(\alpha,\gamma\ge0\), \(C_m,C_o,c_o,c_B,\underline p>0\), \(c_B\le C_m\), \(c_B\le C_o\), \(0<\underline p\le 1/4\), and \(8c_B<\log 5\).  In addition, when \(\gamma>0\) and \(\alpha>0\), assume \(c_B\le C_o c_o^{-\alpha/\gamma}\), and when \(\gamma>0\) and \(\alpha=0\), assume \(c_B\le C_o 4^{-1/\gamma}\).  If the policy class \(\Pi\) is nonempty and every \(\pi\in\Pi\) is measurable, then for the minimax regret \(M_n(\alpha,\gamma)\) of \cref{obj:def:minimax-regret} over the law class \(\mathcal P_{\alpha,\gamma}\) of \cref{obj:def:law-class}, there is a constant \(c>0\), independent of \(n\), such that for all sufficiently large \(n\),
\[
M_n(\alpha,\gamma)
=
\inf_{\widehat\pi}
\sup_{P\in\mathcal P_{\alpha,\gamma}}
\mathbb E_P R_P(\widehat\pi)
\ge
c\,n^{-r_\star(\alpha,\gamma)}.
\]
\end{theoremv}

\cref{obj:thm:minimax-lower} is an observed-law lower bound stated entirely in terms of the law class in \cref{obj:def:law-class}, the regret functional in \cref{obj:def:welfare-regret}, and the minimax risk in \cref{obj:def:minimax-regret}.  \Cref{obj:oeq:feasible-upper} develops the complementary procedure-specific achievability analysis.

Writing the exponent of \cref{obj:def:exponents} explicitly, \(r_\star(\alpha,\gamma)=(1+\alpha)/(2+\alpha+\beta_{\alpha,\gamma})\).  Under strict overlap, \(\beta_{\alpha,\gamma}=0\), and the lower-bound exponent reduces to the usual margin-driven benchmark.  When \(\gamma>0\), the additional term \(\beta_{\alpha,\gamma}\) reflects the loss of treatment-arm information in the small-contrast region.  The section that follows considers a clipped cross-fitted AIPW empirical welfare rule and gives a separate conditional upper bound under explicit nuisance-rate and empirical-process side conditions.
